\begin{thebibliography}{50}

\bibitem{goemotions}
Dorottya Demszky, Dana Movshovitz-Attias, Jeongwoo Ko, Alan Cowen, Gaurav Nemade, and Sujith Ravi.
\newblock GoEmotions: A dataset of fine-grained emotions.
\newblock In \textit{Proceedings of the 58th Annual Meeting of the Association for Computational Linguistics}, pages 4040--4054, 2020.
\newblock \url{https://research.google/pubs/goemotions-a-dataset-for-fine-grained-emotion-classification/}

\bibitem{bert}
Jacob Devlin, Ming-Wei Chang, Kenton Lee, and Kristina Toutanova.
\newblock BERT: Pretraining of deep bidirectional transformers for language understanding.
\newblock \textit{arXiv preprint arXiv:1810.04805}, 2018.

\bibitem{tda}
Herbert Edelsbrunner, David Letscher, and Afra Zomorodian.
\newblock Topological persistence and simplification.
\newblock \textit{Discrete and Computational Geometry}, 28(4):511--533, 2002.

\bibitem{giotto_tda}
Guillaume Tauzin, Umberto Lupo, Lewis Tunstall, Julian B. Rose, Matteo Caorsi, Anibal M. Medina-Mardones, Alberto Becca, and Christian Otani.
\newblock giotto-tda: A Topological Data Analysis Toolkit for Machine Learning and Data Exploration.
\newblock \textit{Journal of Machine Learning Research}, 22(39):1--6, 2021.
\newblock \url{https://jmlr.org/papers/v22/20-1174.html}

\bibitem{tda_ambiguity}
A. V. Smith et al.
\newblock Topological quantification of ambiguity in semantic search.
\newblock \textit{arXiv preprint arXiv:1901.08458}, 2019.
\newblock \url{https://arxiv.org/abs/1901.08458}

\bibitem{mdpi_annotation}
An Experimental Analysis of Data Annotation Methodologies for Emotion Detection in Short Text Posted on Social Media.
\newblock \textit{MDPI}, 2023.
\newblock \url{https://www.mdpi.com/journal/electronics}

\bibitem{affective_tda}
Shaun Canavan.
\newblock AffectiveTDA: Using Topological Data Analysis to Improve Analysis and Explainability in Affective Computing.
\newblock \url{https://scanavan.github.io/AffectiveTDA/}

\bibitem{bert_emotion}
Emotion Classification Using BERT: A Comprehensive Study.
\newblock \textit{Journal of Propulsion Technology}, 2024.
\newblock \url{http://propulsiontechjournal.com/}

\bibitem{blowfish}
Blowfish: Topological and statistical signatures for quantifying ambiguity in semantic search.
\newblock \textit{arXiv}, 2021.
\newblock \url{https://arxiv.org/abs/2105.14371}

\bibitem{tda_review}
Topological data analysis and topological deep learning beyond persistent homology: a review.
\newblock \textit{PMC}, 2023.
\newblock \url{https://pmc.ncbi.nlm.nih.gov/articles/PMC10164632/}

\bibitem{tdavec}
TDAvec: Computing Vector Summaries of Persistence Diagrams for Topological Data Analysis in R and Python.
\newblock \textit{arXiv}, 2022.
\newblock \url{https://arxiv.org/abs/2211.05063}

\bibitem{llm_goemotions}
Large Language Models on Fine-grained Emotion Detection Dataset with Data Augmentation and Transfer Learning.
\newblock \textit{arXiv}, 2023.
\newblock \url{https://arxiv.org/abs/2307.02706}

\bibitem{goemotions_comparison}
Fine-Grained Emotion Detection on GoEmotions: Experimental Comparison of Classical Machine Learning, BiLSTM, and Transformer Models.
\newblock \textit{arXiv}, 2021.
\newblock \url{https://arxiv.org/abs/2105.02496}

\bibitem{emotion_topology}
Emotion topology: extracting fundamental components of emotions from text using word embeddings.
\newblock \textit{PMC}, 2021.
\newblock \url{https://pmc.ncbi.nlm.nih.gov/articles/PMC8323632/}

\bibitem{ai_psychology}
AI meets psychology: an exploratory study of large language models' competence in psychotherapy contexts.
\newblock \textit{Taylor \& Francis}, 2023.
\newblock \url{https://www.tandfonline.com/doi/full/10.1080/09515089.2023.2166542}

\bibitem{early_detection_twitter}
A Multi-Class Deep Learning Approach for Early Detection of Depressive and Anxiety Disorders Using Twitter Data.
\newblock \textit{Unisalento}, 2022.
\newblock \url{http://airus.unisalento.it/}

\bibitem{modernbert_goemotions}
ModernBERT-base-go-emotions.
\newblock \textit{Hugging Face}, 2024.
\newblock \url{https://huggingface.co/cirimus/modernbert-base-go-emotions}

\bibitem{tda_eeg_tool}
Topological Data Analysis as a New Tool for EEG Processing.
\newblock \textit{PMC}, 2021.
\newblock \url{https://pmc.ncbi.nlm.nih.gov/articles/PMC7961298/}

\bibitem{tda_cardio}
Topological Data Analysis in Cardiovascular Signals: An Overview.
\newblock \textit{PMC - NIH}, 2022.
\newblock \url{https://pmc.ncbi.nlm.nih.gov/articles/PMC9317584/}

\bibitem{sarcasm_senti}
Sentiment Analysis: Why Your Model Misses Sarcasm.
\newblock \textit{Label Your Data}, 2023.
\newblock \url{https://labelyourdata.com/articles/sentiment-analysis-sarcasm}

\bibitem{persistence_kernels}
Persistence kernels for classification: A comparative study.
\newblock \textit{Math-Unipd}, 2019.
\newblock \url{http://www.math.unipd.it/}

\bibitem{teaspoon}
Teaspoon: A Python Package for Topological Signal Processing.
\newblock \textit{The Journal of Open Source Software}, 2020.
\newblock \url{https://theoj.org/joss-papers/joss.02507/10.21105.joss.02507.pdf}

\bibitem{persistent_laplacian}
The Persistent Laplacian for Data Science: Evaluating Higher-Order Persistent Spectral Representations of Data.
\newblock \textit{Proceedings of Machine Learning Research}, 2020.
\newblock \url{https://proceedings.mlr.press/v145/memoli21a.html}

\bibitem{giotto_docs}
VietorisRipsPersistence documentation.
\newblock \textit{giotto-tda}.
\newblock \url{https://giotto-ai.github.io/giotto-tda/0.3.0/modules/generated/gtda.homology.VietorisRipsPersistence.html}

\bibitem{simplicial_smote}
Simplicial SMOTE: Oversampling Solution to the Imbalanced Learning Problem.
\newblock \textit{ResearchGate}, 2021.
\newblock \url{https://www.researchgate.net/publication/351608665_Simplicial_SMOTE_Oversampling_Solution_to_the_Imbalanced_Learning_Problem}

\bibitem{ts_smote}
TS-SMOTE: An Improved SMOTE Method Based on Symmetric Triangle Scoring Mechanism for Solving Class-Imbalanced Problems.
\newblock \textit{MDPI}, 2022.
\newblock \url{https://www.mdpi.com/2227-7390/10/21/4043}

\bibitem{march_ambiguity}
MARCH: Evaluating the Intersection of Ambiguity Interpretation and Multi-hop Inference.
\newblock \textit{arXiv}, 2023.
\newblock \url{https://arxiv.org/abs/2305.13264}

\bibitem{ludwig_grammar}
Am feeling better | Meaning, Grammar Guide \& Usage Examples.
\newblock \textit{Ludwig.guru}.
\newblock \url{https://ludwig.guru/s/am+feeling+better}

\bibitem{xgboost_eeg}
Multi-Class Emotion Classification with XGBoost Model Using Wearable EEG Headband Data.
\newblock \textit{SMU Scholar}, 2020.
\newblock \url{https://scholar.smu.edu/datascience_analytics_etds/12/}

\bibitem{tda_info_theory}
An Information Theory of Persistent Homology: Entropy, the Data Processing Inequality, and Rate-Distortion Bounds for Topological Features.
\newblock \textit{MDPI}, 2022.
\newblock \url{https://www.mdpi.com/1099-4300/24/11/1655}

\bibitem{michel}
Paul Michel et al.
\newblock Does the geometry of word embeddings help document classification?
\newblock In \textit{Proceedings of the 2nd Workshop on Representation Learning for NLP}, 2017.

\bibitem{sbert}
Nils Reimers and Iryna Gurevych.
\newblock Sentence-BERT: Sentence embeddings using siamese BERT-networks.
\newblock In \textit{Proceedings of EMNLP}, 2019.

\bibitem{zhu}
Xiaojin Zhu.
\newblock Persistent homology: An introduction and a new text representation for NLP.
\newblock In \textit{Proceedings of IJCAI}, 2013.

\end{thebibliography}
